% Regression: hard05_OL — renders O_L(M) as a vertical chain of M boxes with dotted wire, top trace, and bottom cup.
% Formula: T05 — RMP MPDO-O_L: O_L(M) = a VERTICAL chain of M boxes joined by a
% T05 — RMP MPDO-O_L: O_L(M) = a VERTICAL chain of M boxes joined by a
% dotted wire, trace-closed at the top by an arc through the top box and
% cup-closed at the bottom: an environment with a big trace loop through
% a box, in vertical orientation.
\documentclass[varwidth=19cm, border=6pt]{standalone}
\usepackage{amsmath, amssymb}
\usepackage{tenkz}
\begin{document}

(a) grid attempt — rotate the figure 90 degrees to fit the grid's
horizontal-only chains (periodic closes rows to themselves, one arc is
spurious; boxes carry the vertical wire as physical pairing, which does
not exist for a single wire row):
\[
  O_L(M) \;=\;
  \begin{tenkz}[periodic]
    \tn[box]{M} & \tndots & \tn[box]{M}
  \end{tenkz}
\]

(b) free-tier attempt — raw coordinates, vertical layout, hand-built
racetrack:
\[
  O_L(M) \;=\;
  \begin{tenkzfree}
    \tnput[box]{m1}{(0,0.9)}{M}
    \tnput[box]{m2}{(0,-0.9)}{M}
    \tnjoin{m1.south}{0,0.3}
    \tnjoin{m2.north}{0,-0.3}
    \tnjoin[route=curve, out=90, in=90]{m1.west}{m1.east}
    \tnjoin[route=curve, out=270, in=270]{m2.west}{m2.east}
    \tnjoin{m1.west}{-0.75,0.9}
    \tnjoin{m2.west}{-0.75,-0.9}
    \tnjoin{m1.east}{0.75,0.9}
    \tnjoin{m2.east}{0.75,-0.9}
  \end{tenkzfree}
\]
(vertical ellipsis between the boxes is not expressible: no dots glyph in the
free tier; the top trace arc should pass OVER the box from its left wire to
its right wire, outside the glyph)
\end{document}
